\documentclass[11pt]{article}
\usepackage{fontspec,amsmath,amssymb,amsthm,hyperref}
\usepackage[margin=1in]{geometry}
\setlength{\emergencystretch}{3em}
\title{Pointwise recursion and the L\ensuremath{{}^{2}} algorithm law}
\author{LangevinLean mathematical exposition}
\date{}
\begin{document}\maketitle
\noindent Source module: \texttt{IPLA/Dynamics.lean}. Compile this proof document with XeLaTeX or LuaLaTeX. The accompanying HTML pairs the mathematical argument with the exact Lean statements and records the scope of the formalization.


The formal convergence proof uses equivalence classes in \(L^2\). This module proves that their pushforward laws equal those of the displayed coordinate algorithm, including the exact square-root noise coefficients.


\section{Recover the original pairs from scaled coordinates}


For every particle \(i\), checking parameter and latent coordinates gives \(A_i z=\operatorname{pair}(p_\theta z,\sqrt N\,\operatorname{block}_i z)\). Substitution into the finite sum proves

\[V_N(z)=\sum_iU\bigl(p_\theta z,\sqrt N\,\operatorname{block}_iz\bigr).\]

This identifies the abstract lifted potential with the sum of the original pair potentials at the physical particles.




\textbf{Correspondence with the displayed declarations.}


\begin{itemize}

\item \texttt{joint\_eq\_pair}: A joint extraction is exactly the original parameter-particle pair after unscaling its particle.
\item \texttt{potential\_eq\_pair\_sum}: The lifted potential is the sum of the original pair potential, with physical particles equal to \texttt{sqrt N} times the scaled particle blocks.
\end{itemize}


\noindent Lean declarations: \texttt{joint\_eq\_pair}, \texttt{potential\_eq\_pair\_sum}.

\section{The explicit recursion and its two block formulas}


Set \(Z_0=z_0\) and recursively let \(Z_{k+1}=\operatorname{scaledUpdate}(U,\gamma,Z_k,\xi_k)\). The zero and successor identities are the two defining cases of this recursion. Applying the coordinate identities from the geometry module at \(Z_k\) proves

\[\theta_{k+1}=\theta_k-\frac\gamma N\sum_i\nabla_\theta U(\theta_k,x_k^i)+\sqrt{2\gamma/N}\,\xi_k^\theta,\]\[x_{k+1}^i=x_k^i-\gamma\nabla_xU(\theta_k,x_k^i)+\sqrt{2\gamma}\,\xi_k^i,\qquad x_k^i=\sqrt N(Z_k)_i.\]

The particle formula requires \(N>0\), exactly where unscaling cancels \(N\).




\textbf{Correspondence with the displayed declarations.}


\begin{itemize}

\item \texttt{scaledChain}: IPLA as an ordinary recursion on scaled states and noise vectors.
\item \texttt{scaledChain\_parameter}: The parameter recurrence in the pointwise IPLA chain.
\item \texttt{scaledChain\_particle}: The physical-particle recurrence in the pointwise IPLA chain.
\end{itemize}


\noindent Lean declarations: \texttt{scaledChain}, \texttt{scaledChain\_zero}, \texttt{scaledChain\_succ}, \texttt{scaledChain\_parameter}, \texttt{scaledChain\_particle}.

\section{Induction identifies the L\ensuremath{{}^{2}} and pointwise algorithms}


At time zero, the representative of the \(L^2\) chain is the initialization. Suppose its representative equals the pointwise chain almost surely at time \(k\). The shared update representative theorem gives the Euler formula almost surely; intersect that full-measure event with the induction event, substitute the previous-state equality, and rewrite \(2(\gamma/N)=2\gamma/N\). This proves almost-sure equality at \(k+1\). Finally, almost-everywhere equal measurable maps have equal pushforward measures, hence

\[\rho_k=P\circ\bigl[\omega\mapsto\operatorname{scaledChain}(U,\gamma,Z_0(\omega),\xi(\omega),k)\bigr]^{-1}.\]


The representative induction uses equality almost everywhere at each fixed finite time. In the successor case the intersection of the update's representative event and the previous induction event still has full measure, so substitution is legitimate. Once the random variables agree almost everywhere, \texttt{Measure.map\_congr} gives equality of their laws. This step is what permits results proved for the Hilbert-space recursion to be presented as results for the explicit coordinate algorithm displayed above.




\textbf{Correspondence with the displayed declarations.}


\begin{itemize}

\item \texttt{lpChain\_coe}: The L\ensuremath{{}^{2}} chain represents the explicit IPLA recursion almost everywhere. This includes the exact identity between \texttt{sqrt (2 * (gamma/N))} and the noise coefficient \texttt{sqrt (2*gamma/N)} in the displayed IPLA update.
\item \texttt{laws\_eq\_scaledChain}: The laws appearing in the final theorem are exactly the laws of the explicit, noise-driven, pointwise IPLA recursion.
\end{itemize}


\noindent Lean declarations: \texttt{lpChain\_coe}, \texttt{laws\_eq\_scaledChain}.

\end{document}
